\documentclass[12pt, a4paper]{article}

\usepackage[margin=2.5cm]{geometry}
\usepackage[utf8]{inputenc}
\usepackage[T1]{fontenc}
\usepackage{microtype}
\usepackage{hyperref}
\usepackage[style=apa, backend=biber]{biblatex}

\addbibresource{references.bib}

\title{Coach-Player Paradigm in Multi-Agent Systems:\\
A Literature Review}
\author{}
\date{}

\begin{document}

\maketitle

\section*{Literature Review}

The coach-player paradigm formalizes a persistent tension in multi-agent systems: how to exploit
global information for coordination while preserving decentralized execution. The foundational
theoretical substrate is the \textit{Centralized Training with Decentralized Execution} (CTDE)
framework, instantiated by MADDPG \parencite{lowe2017maddpg} and QMIX
\parencite{rashid2018qmix}, where a centralized critic with access to joint observations trains
agents that later act from local observations alone. CTDE is effectively a training-time
coach---it holds global knowledge but dissolves before deployment. The most explicit
operationalization of a live, runtime coach is \textbf{COPA} \parencite{liu2021copa}, which
introduces a persistent coach agent with full environmental observability that distributes
individualized strategies to players operating under partial views. COPA couples attention
mechanisms for selective information routing, a variational regularization objective to prevent
strategy collapse, and an adaptive communication protocol that sustains high performance while
the coach intervenes as little as 13\% of the time---demonstrating that sparse, asymmetric
guidance can match or exceed full-information baselines. Critically, COPA generalizes zero-shot
to unseen team compositions across StarCraft, resource collection, and rescue benchmarks,
establishing that coach-player role separation confers structural robustness beyond what flat
peer-coordination designs (e.g., QMIX or CommNet) can achieve when team membership is dynamic.

In LLM-based multi-agent systems, the coach-player paradigm resurfaces through a structurally
analogous but mechanistically distinct lens: \textit{role specialization via prompting}.
\textbf{CAMEL} \parencite{li2023camel} instantiates a three-role architecture---a task specifier
(coach) that refines and constrains the objective, an AI assistant (player) that executes, and an
AI user that drives task progression---with inception prompting ensuring agents remain aligned
with the task specification without human re-intervention. \textbf{MetaGPT}
\parencite{hong2024metagpt} extends this to software engineering workflows by encoding
Standardized Operating Procedures (SOPs) that assign Product Manager and Architect agents as
planning-layer coaches over Engineer and Tester agents as executing players; structured
message-passing and role-gated subscriptions replace the attention-based communication channel
of COPA, yielding a more rigid but scalable coordination fabric. \textbf{AutoGen}
\parencite{wu2024autogen} generalizes further with a programmable orchestrator-worker topology,
where a \texttt{GroupChatManager} or \texttt{SelectorGroupChat} dynamically routes tasks to
specialized agents---an implicit coaching layer operating purely through conversational steering
without fine-tuning. Across these systems, a common finding emerges: explicit role
asymmetry---whether implemented as a trained coach agent \parencite{liu2021copa}, a prompted
task specifier \parencite{li2023camel}, or an orchestration layer
\parencite{hong2024metagpt, wu2024autogen}---consistently reduces error compounding, improves
task coherence, and enables generalization that fully peer-symmetric or single-agent designs
struggle to match.

\printbibliography

\end{document}
